% !TeX encoding = UTF-8
\documentclass[variant=docente,cover=institutional,mode=screen]{ua-engineering-v6-article}
% !TeX encoding = UTF-8
\UASetUniversity{Universidad Austral}
\UASetFaculty{Facultad de Ingeniería}
\UASetDepartment{Departamento de Computación}
\UASetCampus{Pilar}
\UASetCareer{Ingeniería Informática}
\UASetCommission{Comisión A}
\UASetProfessor{Equipo docente}
\UASetSemester{Segundo cuatrimestre 2026}
\UASetCourse{Sistemas Distribuidos}
\UASetDate{2026}


\UASetClassNumber{15}
\UASetClassTitle{Diseño integral de sistemas distribuidos}
\UASetDocKind{Guía docente}

\title{Clase 15 --- Guía docente}
\author{\UAProfessor}
\date{\UADate}

\begin{document}
\maketitle

\section{Objetivo pedagógico profundo}

La meta de esta clase es cerrar el ciclo del curso: pasar desde conocimientos fragmentados sobre mecanismos distribuidos hacia la capacidad de integrarlos en una propuesta coherente y defenderla técnicamente.

La idea que debe quedar instalada es:

\begin{center}
\textbf{un sistema distribuido bien diseñado no es una colección de patrones, sino una argumentación coherente bajo restricciones reales}
\end{center}

\section{Resultados de aprendizaje esperados}

Al finalizar, el estudiante debería poder:

\begin{itemize}
\item usar el esquema $D=(P,F,G,S,T)$ como marco de defensa;
\item argumentar una arquitectura completa sin caer en slogans;
\item justificar trade-offs frente a objeciones;
\item comparar casos de dominio sin aplicar una receta única;
\item reconocer cuándo una propuesta es más marketing que diseño técnico.
\end{itemize}

\section{Estructura sugerida para 90 minutos}

\subsection*{Bloque 1 (15 min): cambiar el nivel de análisis}
Abrir con la idea de que el curso ya no está en “qué hace Raft” o “qué es un SLO”, sino en cómo todos esos mecanismos se integran.

\paragraph{Objetivo didáctico}
Hacer visible el cambio de escala: del mecanismo al sistema completo.

\subsection*{Bloque 2 (20 min): marco de defensa}
Presentar explícitamente el esquema:
\[
D = (P,F,G,S,T)
\]

y usarlo para ordenar qué debe contener una propuesta seria.

\paragraph{Objetivo didáctico}
Dar una estructura reusable para diseño y defensa.

\subsection*{Bloque 3 (20 min): trade-offs}
Trabajar comparaciones:
\begin{itemize}
\item consistencia fuerte vs eventual;
\item monolito vs microservicios;
\item RPC vs eventos;
\item transacción fuerte vs saga.
\end{itemize}

\paragraph{Objetivo didáctico}
Que el alumno aprenda a argumentar decisiones, no a repetir preferencias personales.

\subsection*{Bloque 4 (20 min): revisión de casos}
Tomar tres dominios distintos:
\begin{itemize}
\item e-commerce;
\item finanzas;
\item streaming.
\end{itemize}

\paragraph{Objetivo didáctico}
Mostrar que no hay arquitectura universalmente correcta.

\subsection*{Bloque 5 (15 min): objeciones y defensa}
Simular preguntas hostiles:
\begin{itemize}
\item ¿por qué no un monolito?;
\item ¿qué pasa bajo partición?;
\item ¿por qué no consistencia fuerte total?;
\item ¿cómo operás esto?
\end{itemize}

\paragraph{Objetivo didáctico}
Preparar al estudiante para defensa técnica real, no solo para exposición expositiva.


\section{Qué explicar y por qué}

\subsection*{1. Marco $D=(P,F,G,S,T)$}
Porque organiza todo el razonamiento del curso en una forma defendible.

\subsection*{2. Trade-offs}
Porque la defensa arquitectónica se derrumba si el alumno cree que todo puede maximizarse al mismo tiempo.

\subsection*{3. Casos}
Porque muestran que la arquitectura correcta depende del dominio y no de una receta genérica.

\subsection*{4. Objeciones}
Porque obligan a transformar conocimiento en argumentación técnica.

\section{Preguntas disparadoras recomendadas}

\begin{itemize}
\item ¿Qué problema concreto resuelve esta complejidad?
\item ¿Qué falla es más costosa en este dominio?
\item ¿Qué parte exige consistencia fuerte y cuál no?
\item ¿Qué costo operacional trae esta decisión?
\item ¿Qué parte del sistema no está realmente defendida todavía?
\end{itemize}

\section{Errores frecuentes del alumnado}

\begin{itemize}
\item presentar tecnologías sin problema bien formulado;
\item esconder trade-offs;
\item defender microservicios como símbolo de madurez;
\item suponer consistencia fuerte global sin costo;
\item no incluir operación, observabilidad o incidentes en la propuesta;
\item responder objeciones con slogans en vez de supuestos y garantías.
\end{itemize}

\section{Ejercicio en vivo sugerido}

Dividir la clase en grupos y asignarles:
\begin{itemize}
\item e-commerce global;
\item ledger financiero;
\item pipeline de streaming.
\end{itemize}

Pedir que preparen una mini-defensa usando:
\[
D = (P,F,G,S,T)
\]

Luego someterlos a dos objeciones hostiles por grupo.

\paragraph{Objetivo}
Que el alumno practique la transición desde diseño descriptivo hacia defensa argumentada.

\section{Momento crítico de la clase}

El punto más importante aparece cuando el estudiante entiende que:

\begin{center}
\textbf{la calidad de una arquitectura distribuida se mide por su coherencia defendible bajo objeciones, no por la cantidad de componentes que contiene}
\end{center}

Ese momento debe remarcarse, porque resume el espíritu del cierre del curso.

\section{Criterio de éxito docente}

La clase salió bien si el estudiante puede:

\begin{itemize}
\item estructurar una propuesta con el marco del curso;
\item responder objeciones sin improvisación vacía;
\item reconocer costos explícitos de su diseño;
\item distinguir entre propuesta sólida y propuesta decorativa.
\end{itemize}

\section{Conexión con la clase siguiente}

Esta clase prepara perfectamente la clase final, donde el estudiante ya no solo diseña, sino que \emph{defiende} un sistema ante un jurado o comité técnico.

\begin{uakeybox}
Pregunta puente: si ya podemos integrar y argumentar una arquitectura, ¿cómo se organiza una defensa final convincente ante evaluación crítica y preguntas de comité?
\end{uakeybox}


\section{Plan de conducción ampliado}
La clase debe alternar explicación, predicción y contraste. Antes de revelar cada mecanismo, pedir al grupo que anticipe qué puede salir mal; después, volver al mismo escenario y preguntar qué propiedad mejora y qué costo aparece. Cada bloque conceptual debe cerrar con una decisión de diseño observable.
\subsection*{Secuencia recomendada}
\begin{enumerate}
\item Recuperar el prerequisito de la clase anterior.
\item Presentar un caso mínimo que falle y solicitar hipótesis.
\item Formalizar el concepto central de Diseño integral de sistemas distribuidos, distinguiendo seguridad, progreso y supuestos.
\item Resolver un ejemplo completo en pizarra, incluyendo una falla intermedia.
\item Comparar una alternativa de diseño y explicitar trade-offs.
\item Cerrar con una pregunta del Trabajo Final y una pregunta de defensa oral.
\end{enumerate}
\section{Diagnóstico de comprensión durante la clase}
No preguntar solamente ``¿se entendió?''. Usar preguntas discriminantes: ``¿qué cambia si el mensaje llega dos veces?'', ``¿qué propiedad sigue siendo cierta si cae un nodo?'', ``¿esto demuestra seguridad o progreso?'', ``¿qué observa un cliente externo?''. Si el estudiante responde solo con nombres de patrones, pedir una ejecución concreta.
\section{Criterios para corregir ejercicios}
Una respuesta excelente declara supuestos, construye una secuencia verificable, nombra la garantía con precisión, reconoce un costo y conecta la decisión con el dominio. Penalizar afirmaciones absolutas sin condiciones.
\section{Vinculación con el Trabajo Final Integrador}
Reservar entre 5 y 10 minutos para que cada grupo registre una decisión con: problema, alternativa elegida, alternativa descartada y razón. Esto crea trazabilidad entre las clases y las entregas acumulativas.
\section{Lectura docente previa}
Revisar la bibliografía indicada en los apuntes y preparar al menos un contraejemplo que no aparezca literalmente en las diapositivas. El objetivo es poder responder preguntas de frontera sin convertir la clase en una lectura del deck.

\section{Hito del Trabajo Final: consolidación}
La especificación recibida ubica en torno a esta clase la consolidación de arquitectura y el prototipo. Antes de evaluar, el equipo docente debe confirmar el calendario definitivo, porque el documento fuente contiene referencias distintas para la entrega de Etapa 3 y la defensa final. La evaluación técnica debe seguir la versión de calendario comunicada formalmente a los estudiantes.

\section{Arquitectura didáctica v9 para 180 minutos}
\subsection*{Bloque 1 - desafío inicial (20 min)}
Presentar una situación donde aparezca supuestos ocultos, garantías incompatibles, operación ausente y complejidad injustificada. Pedir predicciones individuales antes de introducir vocabulario. El objetivo es hacer visibles los supuestos intuitivos y recuperar el prerequisito de la clase anterior.

\subsection*{Bloque 2 - construcción conceptual (35 min)}
Formalizar architecture review, ADRs, trade-offs, modelo de fallas, evidencia y defensa de decisiones. Separar seguridad, progreso y experiencia del usuario. Explicar no solo \emph{qué} hace cada mecanismo, sino \emph{por qué} existe y qué supuesto permite que funcione.

\subsection*{Bloque 3 - modelo y figura (35 min)}
Construir en pizarra una línea temporal, grafo, log, máquina de estados o tabla de decisiones. Recorrer una ejecución nominal y una adversarial. La representación debe quedar como referencia para los apuntes y el ejercicio principal.

\subsection*{Pausa (10 min)}
Recoger una pregunta anónima o un punto de confusión. Elegir una pregunta que permita distinguir síntoma, causa y propiedad.

\subsection*{Bloque 4 - caso trabajado con falla (40 min)}
Aplicar el mecanismo al caso conductor. Introducir una falla a mitad de ejecución y detenerse en cada transición: quién sabe qué, qué se persiste, qué garantía sigue vigente y qué observa el cliente.

\subsection*{Bloque 5 - taller de decisión (25 min)}
Los grupos aplican P-F-G-S-T a su Trabajo Final. Deben producir un ADR breve con alternativa elegida, alternativa descartada y evidencia pendiente. La evidencia de esta clase es arquitectura, alternativas, ADRs, experimento y defensa ante objeciones.

\subsection*{Bloque 6 - cierre y exit ticket (15 min)}
Cada estudiante responde: propiedad principal, supuesto decisivo, costo aceptado y condición bajo la cual cambiaría de diseño. Usar las respuestas para iniciar la clase siguiente.

\section{Qué explicar, por qué y cómo verificarlo}
\begin{itemize}
\item Explicar el problema antes del producto, porque reconocer la familia de problema es el resultado cognitivo central.
\item Explicar el invariante, porque permite evaluar configuraciones y no solo repetir un flujo nominal.
\item Explicar el límite, porque supuestos ocultos, garantías incompatibles, operación ausente y complejidad injustificada puede invalidar supuestos o reducir progreso.
\item Verificar comprensión mediante arquitectura, alternativas, ADRs, experimento y defensa ante objeciones, no mediante una definición aislada.
\end{itemize}

\section{Preguntas socráticas}
\begin{itemize}
\item ¿Qué sabe cada actor en este punto de la ejecución?
\item ¿Qué propiedad sigue siendo cierta si la respuesta nunca llega?
\item ¿Qué parte es safety y cuál es liveness?
\item ¿Qué alternativa más simple resolvería el problema si relajamos una garantía?
\item ¿Qué observación refutaría nuestra hipótesis de diseño?
\end{itemize}

\section{Errores previsibles y estrategia de corrección}
\begin{itemize}
\item Si el alumno responde con un nombre de tecnología, pedir actores, mensajes y estados.
\item Si confunde timeout con falla, construir dos ejecuciones indistinguibles.
\item Si promete todas las garantías, pedir comportamiento bajo partición o saturación.
\item Si mide solo rendimiento, preguntar cómo evidencia la propiedad semántica.
\item Si la solución es excesiva, pedir la alternativa mínima y el costo marginal de la garantía fuerte.
\end{itemize}

\section{Criterios de evaluación formativa}
Una respuesta excelente declara supuestos, recorre una ejecución, nombra con precisión la garantía, reconoce el costo y propone evidencia reproducible. Una respuesta suficiente identifica el mecanismo y el trade-off principal. Una respuesta insuficiente usa slogans, omite fallas o confunde producto con propiedad.

\section{Preparación docente y bibliografía}
Lectura previa recomendada: Architecture Decision Records; DDIA como marco comparativo; revisiones de sistemas reales. Preparar un contraejemplo adicional y una variante del caso en la que cambie el costo del error; esto evita convertir la clase en una lectura de diapositivas.

\end{document}
